\bookchapter{Logging, Replay and Visualisation}\label{ch:ros5}

\section*{Explanation}
A robot that works only while you watch it is a demo, not a robot. At a competition you have few minutes between games, the field is noisy, and the bug that appeared in the last game will not appear again on command. Three tools solve this: \textbf{logging} (the programs tell you what they think and do), \textbf{recording and replay} (\emph{rosbag}: all data of a run is stored and can be played back as often as you like, on any computer), and \textbf{visualisation} (you \emph{see} the frames, the detections, the plans). Together they turn ``it did something strange'' into a reproducible test.

\section*{Logging}
Every node has a logger with five levels: \verb!DEBUG! (details for developers), \verb!INFO! (normal events), \verb!WARN! (something is wrong, but the robot continues), \verb!ERROR! (a function failed) and \verb!FATAL! (the node cannot continue).
\begin{verbatim}
self.get_logger().debug(f"candidates: {len(contours)}")
self.get_logger().info("ball found")
self.get_logger().warn("camera frame is 120 ms old")
self.get_logger().error(f"tf lookup failed: {e}")
self.get_logger().info("searching", throttle_duration_sec=1.0)
self.get_logger().info("kickoff!", once=True)
\end{verbatim}
\verb!throttle_duration_sec! prints at most once per second: without it a message in a $30$\,Hz callback produces $1800$ lines per minute, hides everything else and slows the node down. Messages go to the terminal, to a file in \verb!~/.ros/log! and to the topic \verb!/rosout! (so that tools can show messages of all nodes in one place: \verb!rqt_console!). The level of a node is set at start, without changing the code:
\begin{verbatim}
$ ros2 run my_robot_vision ball_detector --ros-args --log-level debug
$ ros2 run my_robot_vision ball_detector --ros-args \
      --log-level ball_detector:=debug --log-level rcl:=warn
\end{verbatim}
Recent releases also let you choose the logging backend (the environment variable \verb!RCL_LOGGING_IMPLEMENTATION!) and the level per message; see the release notes of your distribution.

\begin{keyidea}[What a good log message contains]
\textbf{Who} (the node is in the log automatically), \textbf{what happened}, and \textbf{the numbers} that you will need later. ``Error'' is useless; ``tf lookup \verb!map!$\leftarrow$\verb!camera_link! failed: extrapolation 0.31\,s into the future'' is a bug report. Log \emph{state changes} (``searching $\to$ approaching ball'') at INFO and never every cycle; log everything that you would like to know afterwards \emph{before} the competition, not after.
\end{keyidea}

\section*{Recording with rosbag2}
A \textbf{bag} is a file (or a folder) with time-stamped messages of selected topics. \verb!ros2 bag record! subscribes to the topics, exactly like any node, and writes what it receives.
\begin{verbatim}
$ ros2 bag record -o match1 /camera/image_compressed /imu/data \
      /joint_states /tf /tf_static /ball /cmd_vel
$ ros2 bag info match1
Files:             match1_0.mcap
Bag size:          1.2 GiB
Storage id:        mcap
Duration:          612.4s
Messages:          453208
Topic information: Topic: /ball | Type: ... | Count: 18372 ...
\end{verbatim}
(The output is typical; the numbers depend on your system.) Rules for recording that avoid pain later:
\begin{itemize}[nosep]
\item \textbf{Record the inputs and the outputs}: sensor data (compressed images, IMU, joint states), \verb!/tf! and \verb!/tf_static!, and the decisions (\verb!/ball!, \verb!/cmd_vel!). Then you can replay the sensors and compare the new decisions to the old ones.
\item \textbf{Do not record everything} (\verb!-a!): raw images fill the disk and slow down the robot. The sizes below show why.
\item \textbf{Use compressed images} and, if needed, compression of the bag (\verb!--compression-mode file --compression-format zstd!).
\item \textbf{Record the QoS}: a late-joining recorder must subscribe with a QoS that matches the publisher (\verb!tf_static! is transient local): current versions of \verb!ros2 bag record! adapt automatically; if a topic stays empty, check \verb!ros2 topic info -v!.
\item \textbf{Split long recordings} (\verb!--max-bag-size 1000000000! splits at $1$\,GB), so that a crash or a full disk loses at most one file. The storage format is SQLite3 or MCAP depending on the release (MCAP is the recommended one: robust against a power cut, and readable by many tools).
\item Write the recording to a fast local disk, and keep a text file next to it: which robot, which software version (\verb!git rev-parse HEAD!), which parameters, what you tested.
\end{itemize}
\begin{center}
\renewcommand{\arraystretch}{1.25}
\begin{tabular}{p{5.2cm}rr}
\textbf{Topic} & \textbf{data rate} & \textbf{per minute}\\ \hline
Raw image $640\times480$ RGB, $30$\,Hz & $27.6$ MB/s & $1659$ MB\\
JPEG image $45$\,kB, $30$\,Hz & $1.35$ MB/s & $81$ MB\\
Joint states, $22$ joints, $500$\,Hz (about $590$ bytes per message) & $0.29$ MB/s & $17.6$ MB\\
\end{tabular}
\end{center}
A $20$-minute match with the raw image is about $33$\,GB, with JPEG $1.6$\,GB plus about $0.35$\,GB for the joints.

\section*{Playing back}
\begin{verbatim}
$ ros2 bag play match1 --clock           # real speed, publishes /clock
$ ros2 bag play match1 -r 0.5            # half speed
$ ros2 bag play match1 --loop --start-offset 120
$ ros2 bag play match1 --topics /imu/data /joint_states /tf
\end{verbatim}
The nodes under test must run with \verb!use_sim_time:=true! (Chapter~\ref{ch:ros3}): they then take the time from the \verb!/clock! that the player publishes, not from the wall clock, which would make every time stamp in the replayed data ``too old''. The typical workflow after a bad game: (1) copy the bag to your laptop; (2) replay the sensor topics (\emph{not} the old results such as \verb!/ball! and \verb!/cmd_vel!) into your new version of the detector; (3) compare its \verb!/ball! with the recorded one; (4) fix, replay, compare again, until the problem is gone. This is also the basis for \emph{regression tests}: a short bag that once triggered a bug becomes a test that the bug must never come back (Chapter~\ref{ch:py11}).

\subsection*{Reading a bag in Python}
Inside ROS~2 the package \verb!rosbag2_py! reads bags:
\begin{verbatim}
import rosbag2_py
from rclpy.serialization import deserialize_message
from rosidl_runtime_py.utilities import get_message

reader = rosbag2_py.SequentialReader()
reader.open(rosbag2_py.StorageOptions(uri="match1", storage_id="mcap"),
            rosbag2_py.ConverterOptions("", ""))
types = {t.name: t.type for t in reader.get_all_topics_and_types()}
stamps = []
while reader.has_next():
    topic, data, t_ns = reader.read_next()
    if topic == "/ball":
        msg = deserialize_message(data, get_message(types[topic]))
        stamps.append(t_ns * 1e-9)
\end{verbatim}
The same data can be read \emph{without} ROS installed: the MCAP format is an open container, with Python packages (\verb!mcap!, \verb!mcap-ros2-support!) that can be installed with pip on any laptop (not tested here). Therefore you can analyse a match recorded on the robot on a computer that has only Python and numpy.

\begin{example}
\textbf{What does \verb!ros2 topic hz! compute?} Take the time stamps of a topic: $\Delta t_i=t_{i+1}-t_i$. The average rate is $1/\overline{\Delta t}$; a healthy $30$\,Hz stream has $\Delta t\approx33.3$\,ms with a small spread. We simulated $60$\,s of a $30$\,Hz camera with $2$\,ms of timing jitter and $3\,\%$ lost frames and computed the statistics from the time stamps with \verb!numpy.diff!: $1744$ of $1800$ frames arrived, the mean rate was $29.07$\,Hz, the intervals ranged from $24.0$ to $101.6$\,ms (a gap of three periods happens when two frames in a row are lost) with a standard deviation of $6.9$\,ms, and $52$ gaps were longer than $1.5$ periods. The mean alone ($29$\,Hz) looks harmless; the \emph{largest gap} (about $100$\,ms) is what a controller experiences. Always look at the maximum and at the number of gaps, not only at the average.
\end{example}

\section*{Visualisation}
\begin{itemize}[nosep]
\item \textbf{rviz2}: the 3D view of the robot. It shows frames (TF), the robot model (URDF), images, point clouds, maps, paths and \emph{markers}. Set the \textbf{Fixed Frame} first (usually \verb!map!); a display that stays red with ``no transform'' shows a frame problem (Chapter~\ref{ch:ros4}).
\item \textbf{rqt}: small plug-ins in one window: \verb!rqt_graph! (the node/topic graph of Figure~\ref{fig:ros1_graph}), \verb!rqt_console! (log messages with filter), \verb!rqt_plot! and \verb!rqt_reconfigure! (change parameters with sliders).
\item \textbf{PlotJuggler} and \textbf{Foxglove}: more powerful tools for time series, joint torques and replay of bags; Foxglove reads MCAP files and can show a robot live through a bridge node. (Both are separate installations.)
\end{itemize}
To show your own result in \verb!rviz2!, publish a \emph{marker}, a message that says ``draw a sphere here'':
\begin{verbatim}
from visualization_msgs.msg import Marker
m = Marker()
m.header.frame_id = "map"
m.header.stamp = self.get_clock().now().to_msg()
m.ns, m.id = "ball", 0
m.type, m.action = Marker.SPHERE, Marker.ADD
m.pose.position.x, m.pose.position.y = bx, by
m.pose.orientation.w = 1.0
m.scale.x = m.scale.y = m.scale.z = 0.13          # ball diameter
m.color.r, m.color.g, m.color.a = 1.0, 0.5, 1.0   # orange, opaque
m.lifetime.sec = 1                                # disappears if not renewed
self.marker_pub.publish(m)
\end{verbatim}
Do not forget \verb!color.a! (the transparency): the default $0$ makes the marker invisible, the most common ``my marker does not appear''. The \verb!lifetime! makes a stale marker vanish.

\subsection*{Measuring delay}
For every message with a header, the \emph{age} is $\text{now}-\text{stamp}$. A node can compute it per message and log it (throttled), and \verb!ros2 topic delay /camera/image! shows it from the command line. A delay that grows over time means that a subscriber is slower than the publisher and the queue grows (Chapter~\ref{ch:ros1}, problem~10); a constant delay is the sum of capture, transmission and processing time (Chapter~\ref{ch:hw4}).

\begin{aibox}[Bags as training data]
The same bag that helps you debug is a source of training data for neural networks (Part~\ref{part:ai}): images with the detections and the robot's pose give labelled examples from the real robot and real lighting. Record the time stamps and the calibration with them; without them the images cannot be matched to the poses afterwards.
\end{aibox}

\section*{Practice problems}
\begin{enumerate}[label=\textbf{\thechapter.\arabic*}]
  \item A callback at $30$\,Hz logs one INFO line of $80$ bytes per call. How many lines and bytes per hour? Which option avoids this?
  \item Name the log level you would use for (a) ``kickoff signal received'', (b) ``camera image is 300 ms old'', (c) ``cannot open the camera device'', (d) the content of a list of candidate balls.
  \item Estimate the size of a bag of $10$ minutes with JPEG images ($50$\,kB, $30$\,Hz), a $200$\,Hz IMU ($100$ bytes per message) and $500$\,Hz joint states ($600$ bytes).
  \item You replay a bag with \verb!ros2 bag play match1! and your nodes complain that messages are ``too old'' and \verb!tf! lookups fail. What did you forget?
  \item A bag contains the topics \verb!/image!, \verb!/ball! and \verb!/cmd_vel!. You want to test a new detector. Which topics do you replay and which do you let the new nodes produce?
  \item Compute from the time stamps $0.000,\,0.034,\,0.066,\,0.100,\,0.200,\,0.233$\,s the mean rate, the largest gap and the number of gaps that are longer than $1.5$ times the nominal $33.3$\,ms period.
  \item Your rviz marker does not appear. List four causes, ordered from most to least likely.
  \item Why should a recording not be started and stopped by hand in the heat of the game? Describe a way to start it automatically, and how to protect the disk from running full.
  \item What is the advantage of an open container format such as MCAP for the team, compared with a database that only ROS tools can read?
  \item[\textbf{\thechapter.10}$\star$] Write a script that reads the time stamps of one topic from a bag (as in the chapter) and prints the mean rate, the maximum gap, the number of gaps larger than $1.5$ periods and a histogram of the intervals (text or matplotlib). Test the statistics part with the simulated time stamps of the chapter (random generator, $30$\,Hz, $3\,\%$ loss).
\end{enumerate}
